%
% Chapter 1 — Introduction
%
\chapter{Introduction}
\label{chap:intro}
\setlength{\parskip}{1em}

\section{Relevance, Problem Statement, and Research Gap}

Urban transportation management systems rely on heterogeneous data collected from sensors, CSV exports, API feeds, and streaming telemetry~\cite{wang1996beyond,pipino2002data,zanella2014iot}. In practice, these sources often contain missing values, duplicated records, schema inconsistencies, invalid coordinates, delayed timestamps, and domain-specific anomalies. Such defects complicate forecasting, anomaly detection, classification, and other analytical workflows, because the reliability of downstream results depends directly on the quality and reproducibility of the preparation process~\cite{sculley2015hidden}.

Existing studies on automated preprocessing, intelligent transportation systems, and neural traffic analytics usually address these stages separately~\cite{liu2022automated,feurer2019auto,feurer2015autosklearn,its_bigdata2018,nie2023patchtst,yu2018stgcn,zeng2023transformers}. Data cleaning is often treated as a preliminary technical step, while model training, deployment, and monitoring are evaluated independently~\cite{zaharia2018accelerating,kreutzberger2022mlops}. As a result, transportation analysts and system operators may lack a unified mechanism for measuring data quality, applying consistent transformations, preserving preparation history, and connecting prepared datasets to downstream models.

The research problem addressed in this thesis is the lack of an auditable and reproducible workflow for automating data preparation in urban transportation management systems while preserving compatibility with downstream analytical models. This thesis investigates how such a workflow can be designed, implemented, and evaluated through the \textbf{Flowmatic} platform. The dissertation builds on supporting articles developed during the master's programme---on adaptive passenger-flow modelling, automated data preparation, cross-city federated forecasting, and event-driven traffic microservices---from which methodological and empirical insights were collected and applied here~\cite{bakytuly2024hybrid,yedilkhan2025flowmatic,bakytuly2025adaptive,bakytuly2026microservices,sakhipov2026xfedformer}.

\section{Aim and Objectives of the Research}
\label{sec:objectives}

The aim of this research is to develop and evaluate an intelligent software platform for automating data preparation in urban transportation management systems, with support for measurable data quality assessment, reproducible preprocessing, and integration with downstream analytical models.

To achieve this aim, the following objectives are defined:
\begin{enumerate}
  \item To analyze existing approaches to data quality assessment, automated preprocessing, and model lifecycle management in intelligent transportation systems.
  \item To design the architecture of a platform that supports batch and streaming data preparation workflows for urban transportation data.
  \item To implement mechanisms for data quality assessment, cleaning, transformation, storage, and export.
  \item To integrate the preparation workflow with a model registry and inference service for transportation-related analytical tasks.
  \item To evaluate the proposed platform using data quality metrics, latency measurements, controlled corruption experiments, baseline comparisons, and routing validation.
  \item To define the limitations of the proposed approach and identify directions for future deployment in real urban transportation environments.
\end{enumerate}

The object of the research is the data preparation process in urban transportation management systems. The subject of the research is the set of methods and software mechanisms for automated quality assessment, preprocessing, routing, and preparation of transportation data for analytical models.

\section{Scientific Novelty and Practical Significance}

The scientific novelty of the research lies in the integrated evaluation of data preparation, quality assessment, and model-oriented processing within a single reproducible workflow for urban transportation data.

The main elements of novelty are:
\begin{enumerate}
  \item An integrated preparation-to-inference lifecycle that combines data quality assessment, automated cleaning, export, model registry usage, and downstream inference in one workflow.
  \item A six-dimensional Data Quality Index applied to urban transportation data and evaluated under controlled data corruption scenarios~\cite{yedilkhan2025flowmatic}.
  \item A reproducible evaluation protocol linking preparation quality, batch-processing latency, baseline comparisons, neural model performance, and routing validation.
  \item A model-oriented routing mechanism that connects prepared traffic and sensor streams to task-specific analytical models~\cite{bakytuly2026microservices}.
  \item Published artifacts and evaluation outputs that support reproducibility of the experimental results~\cite{kreutzberger2022mlops,wolf2020transformers}.
\end{enumerate}

The practical significance of the research is represented by the implemented Flowmatic platform, which provides authenticated data upload, automated quality assessment, cleaning, storage, export, smart-city simulation workflows, model registry integration, and inference support. The platform can be used as a research prototype for preparing transportation datasets and demonstrating how data preparation can be connected with downstream analytics in urban transportation management systems.

\section{Author's Prior Publications}
\label{sec:author-pubs}

The publications listed below are supporting articles for this master's thesis. While working on these studies, methodological and empirical insights were collected that informed the design, implementation, and evaluation of the Flowmatic platform presented in this dissertation. The principal publications are cited in APA style below and referenced throughout the thesis.

\begin{enumerate}
  \item Begisbayev, D., Sakhipov, A., Seiitbek, R., Mansurova, A., \& Mansurova, A. (2024). Investigation of hybrid models and architectures for real-time adaptive passenger flow prediction. In \emph{2024 7th International Conference on Data Science and Information Technology (DSIT)}, pp.\ 1--5. IEEE. \url{https://doi.org/10.1109/DSIT61374.2024.10880981}
  \item Yedilkhan, D., Sakhipov, A., Seitbek, R., Mektepayeva, A., \& Begisbayev, D. (2025). Intelligent assistant for data preparation automation in urban transportation management systems. \emph{KazATC Bulletin}, 137(2), 310--324. \url{https://doi.org/10.52167/1609-1817-2025-137-2-310-324}
  \item Mektepbayeva, A., Begisbayev, D., Seiitbek, R., Khaimuldin, N., Sakhipov, A., \& Rakhimzhanov, D. (2025). Adaptive machine learning algorithms for data processing in transportation systems. In \emph{2025 IEEE 5th International Conference on Smart Information Systems and Technologies (SIST)}, pp.\ 1--8. IEEE. \url{https://doi.org/10.1109/SIST61657.2025.11139286}
  \item Bakytuly, R. (2025). Real-time passenger flow prediction using automated data preparation. Paper presented at the 10th International Conference on Digital Technologies in Education, Science and Industry (DTESI), 2025.
  \item Sakhipov, A., \& Seiitbek, R. (2026). Event-driven microservices for incident detection and response in intelligent traffic systems. \emph{International Journal of Information and Communication Technologies}, 7(1), 218--230. \url{https://doi.org/10.54309/IJICT.2026.25.1.014}
  \item Sakhipov, A., Uzdenbayev, Z., Begisbayev, D., Mektepbayeva, A., Seiitbek, R., \& Yedilkhan, D. (2026). Deep heterogeneity learning for cross-city transit forecasting: A differentially private federated framework with mixture-of-experts and seasonal decomposition. \emph{Frontiers in Future Transportation}, 7, 1644979. \url{https://doi.org/10.3389/ffutr.2026.1644979}
\end{enumerate}

\noindent The thesis author (Ramazan Seiitbek) is listed as \emph{Seiitbek, R.} on the IEEE, IJICT, and Frontiers items above and as \emph{Bakytuly, R.} on the DTESI 2025 entry, matching each publisher's author metadata.

\noindent A companion manuscript on architectural patterns for Astana smart-traffic incident management (2026, \emph{in preparation}) documents the semi-synthetic data generator used in Appendix~\ref{app:astana}; it is cited as unpublished work~\cite{bakytuly2026architectural} and is not listed above because it is not yet published.


\section{Thesis Structure}

The thesis is organized as follows. Chapter~\ref{chap:literature} reviews related work on data quality, automated preprocessing, intelligent transportation systems, forecasting, anomaly detection, model registries, and reproducibility. Chapter~\ref{chap:methodology} describes the research methodology, including the data quality assessment approach, preparation workflow, routing logic, evaluation protocols, and research questions. Chapter~\ref{chap:implementation} presents the implementation of the Flowmatic platform. Chapter~\ref{chap:results} reports experimental results and discusses preparation quality, latency, baseline comparisons, neural evaluation, and routing validation. Chapter~\ref{chap:conclusion} summarizes the findings, limitations, and future research directions. Appendices document reproducibility commands and Astana dataset construction.
